\section{Performance}
This section provides expected complexity and performance considerations in the
absence of dedicated benchmarks.

\paragraph{Scheduler complexity}
Let $T$ be the number of runnable tasks in an instant, $B$ the number of
blocked tasks, and $S$ the number of registered signals. The scheduler processes
each runnable task at most once per step and checks guard satisfaction by
scanning the task's guard list. Over an instant, the total cost is
$O(T + G)$ where $G$ is the sum of guard-list sizes across tasks, plus
signal finalization at $O(S)$.

\paragraph{Signal operations}
Emitting an event signal resumes all its awaiters immediately (in scheduler
terms), so the cost is proportional to the number of waiters on that signal.
Aggregate signals defer resumption to instant end, so emission cost is $O(1)$
plus the cost of the user-provided combine function. Finalization resumes all
aggregate awaiters once per instant.

\paragraph{Queue sizes and memory}
Memory usage grows with the number of tasks, signals, and registered waiters.
In typical use, each \texttt{await} adds one continuation to a signal's waiter
list, and each guard adds a task to a guard\_waiter list. Worst-case memory
consumption is therefore linear in the number of suspended tasks and guards.

\paragraph{Assumptions}
These bounds assume: (1) guard lists are short; (2) signals are not created at
unbounded rates per instant; (3) the combine function for aggregate signals is
constant-time. Under these assumptions, the scheduler is expected to scale
linearly with the number of active tasks and signals per instant.

\paragraph{Effect and allocation overhead}
Each reactive primitive is implemented as an effect, which captures a
continuation and allocates a task record. The dominant overheads are therefore
continuation capture, task allocation, and queue operations. In practice, this
means performance is sensitive to allocation rates and GC behavior, especially
in programs that spawn many short-lived tasks per instant.

\paragraph{Preemption and guard costs}
Weak preemption introduces a small constant overhead per protected task: a kill
token, a cleanup closure, and a liveness check when the task is enqueued. Guard
checks scale with the number of guard conditions attached to a task, which
makes guard-heavy programs more sensitive to the size of guard lists than to
the number of signals overall.

\paragraph{ReactiveML comparison plan}
To compare \tempo{} against ReactiveML, we propose a small benchmark suite that
captures the common synchronous patterns shared by both systems:
\begin{itemize}
\item \textbf{Ping-pong signals}: two processes alternately emit and await
  signals across instants (measures basic signaling and scheduling overhead).
\item \textbf{Barrier synchronization}: $N$ parallel tasks emit on a barrier
  signal and a coordinator waits for all emissions (measures aggregation and
  guard wakeups).
\item \textbf{Timeout and preemption}: a request handler guarded by a timeout,
  implemented with weak preemption (measures watch/guard overhead).
\item \textbf{Broadcast and fan-out}: one emitter wakes $N$ awaiters in the same
  instant (measures awaiter resumption cost).
\item \textbf{Guarded loops}: repeated guarded activation with pauses across
  many instants (measures long-running scheduling stability).
\end{itemize}
For each scenario, the metrics of interest are throughput (instants per second),
latency per instant, and memory usage as a function of the number of tasks and
signals. These benchmarks are intended as future work; the current draft does
not include measured results.

\paragraph{Benchmark parameters}
\begin{table}[h]
\centering
\begin{tabularx}{\linewidth}{lY}
\hline
Benchmark & Parameters and cost drivers \\ \hline
Ping-pong & $I$ instants; cost dominated by signal wakeups. \\
Barrier   & $N$ tasks; cost dominated by guard wakeups and aggregation. \\
Timeout   & $I$ instants; cost dominated by watch/kill handling. \\
Broadcast & $N$ awaiters; cost dominated by awaiter resumption. \\
Guarded loop & $I$ instants; cost dominated by guard checks per instant. \\ \hline
\end{tabularx}
\caption{Benchmark parameters and expected cost drivers.}
\end{table}

\paragraph{Reproducibility outline}
Each scenario can be implemented as a standalone executable in
\texttt{tests/ok} (\tempo{}) and a corresponding ReactiveML program, compiled with
the ReactiveML toolchain. A minimal protocol should record: (1) number of
instants executed; (2) wall-clock time; (3) task/signal counts. This makes the
comparison transparent even before results are reported.

\paragraph{Concrete protocol}
For each scenario, fix a parameter grid (e.g., number of tasks $N \in
\{10, 10^2, 10^3\}$ and instants $I \in \{10^3, 10^4\}$), and run both \tempo{} and
ReactiveML implementations on the same machine with CPU frequency scaling
disabled. Measure wall-clock time for $I$ instants, record peak resident memory,
and log the number of signals and waiters. Report median over at least 5 runs.

\paragraph{Test coverage map}\label{sec:feature-tests}
The list below shows which \texttt{tests/ok} programs exercise each public
feature. This supports the claim that all features are covered by tests.
\begin{itemize}
\item \textbf{emit/await}:
  \texttt{emit\_once\_await\_one.ml} \\
  \texttt{emit\_once\_await\_many.ml}.
\item \textbf{await\_immediate}: \texttt{emit\_once\_await\_immediate.ml},
  \\ \texttt{when\_await\_imm\_same\_instant.ml},
  \\ \texttt{when\_await\_imm\_next\_instant.ml}.
\item \textbf{pause}: \texttt{pause.ml}.
\item \textbf{parallel}: \texttt{parallel.ml}, \texttt{when\_parallel.ml}.
\item \textbf{guards/when\_}: \texttt{when\_next\_instant.ml},
  \\ \texttt{when\_same\_instant.ml}, \texttt{when\_nested.ml},
  \\ \texttt{when\_complete.ml}.
\item \textbf{watch/weak preemption}: \texttt{watch\_complete.ml},
  \\ \texttt{watch\_nested.ml}, \texttt{watch\_killed.ml},
  \\ \texttt{watch\_when\_complete.ml}, \texttt{watch\_when\_killed.ml}.
\item \textbf{kill}: \texttt{with\_kill\_abort.ml}, \texttt{with\_kill\_inherit.ml}.
\item \textbf{aggregate signals}:
  \texttt{emit\_agg\_collect.ml} \\
  \texttt{emit\_agg\_sum.ml}.
\item \textbf{present\_then\_else}: \texttt{present\_then\_else.ml}.
\end{itemize}

\paragraph{Testing strategy (summary)}
Tests are organized by primitive (signals, guards, preemption) and by
composition patterns (parallel, nested guards, immediate awaits). Each test
has an expected output, enabling regression checks on observable behavior at
instant boundaries.
